\documentclass[aspectratio=169]{beamer}
\usepackage[utf8]{inputenc}
\usepackage[T1]{fontenc}
\usepackage{lmodern}
\usepackage{xcolor}
\usepackage[portuguese]{babel}

\definecolor{BgCream}{HTML}{FAF8F5}
\definecolor{TextEspresso}{HTML}{4A4238}
\definecolor{NudeTaupe}{HTML}{BFAFA6}
\definecolor{SoftBlush}{HTML}{D9C5B2}

\setbeamertemplate{navigation symbols}{}
\setbeamercolor{background canvas}{bg=BgCream}
\setbeamercolor{normal text}{fg=TextEspresso}
\setbeamercolor{structure}{fg=TextEspresso}

\setbeamertemplate{itemize item}{\color{NudeTaupe}\textendash}
\setbeamertemplate{itemize subitem}{\color{NudeTaupe}\(\circ\)}

\setbeamerfont{title}{size=\huge, series=\bfseries}
\setbeamerfont{frametitle}{size=\Large, series=\bfseries}

\setbeamertemplate{title page}{
  \vspace*{1.5cm}
  \begin{center}
    \usebeamerfont{title}\textcolor{TextEspresso}{\inserttitle}\par
    \vspace{0.35cm}
    {\large\textcolor{NudeTaupe}{\insertsubtitle}}\par
    \vspace{0.6cm}
    {\color{NudeTaupe}\rule{2.5cm}{0.5pt}}\par
    \vspace{0.8cm}
    \usebeamerfont{author}\textcolor{TextEspresso}{\insertauthor}\par
    \vspace{0.3cm}
    \usebeamerfont{institute}\small\textcolor{NudeTaupe}{\insertinstitute \quad|\quad \insertdate}
  \end{center}
}

\setbeamertemplate{frametitle}{
  \vspace*{0.8cm}
  \usebeamerfont{frametitle}\insertframetitle\par
  \vspace*{-0.1cm}
  {\color{NudeTaupe}\rule{1.5cm}{0.5pt}}
  \vspace*{0.2cm}
}

\newenvironment{ideiachave}{
  \vspace{0.5cm}
  \begin{center}
  \color{TextEspresso}
  \itshape\Large
}{
  \end{center}
  \vspace{0.5cm}
}

\usepackage{csquotes}

\title{Orquestração de Contentores com Kubernetes}
\subtitle{Infraestruturas de Computação na Nuvem, Aula 05}
\author{Catarina Silva}
\institute{ESTGA, Universidade de Aveiro}
\date{2026}

\begin{document}

\begin{frame}[plain]
  \titlepage
\end{frame}

\begin{frame}
\frametitle{Agenda \textcolor{NudeTaupe}{\small(1/2)}}
\begin{itemize}
    \setlength\itemsep{0.4cm}
    \item Porque precisamos de orquestração além de um único host.
    \item Arquitetura Kubernetes: control plane e worker nodes.
    \item Objetos principais: Pod, Deployment, ReplicaSet, Service, Namespace.
\end{itemize}
\end{frame}

\begin{frame}
\frametitle{Agenda \textcolor{NudeTaupe}{\small(2/2)}}
\begin{itemize}
    \setlength\itemsep{0.4cm}
    \item Ciclo de vida de um Pod e self-healing.
    \item O ciclo de reconciliação e o modelo declarativo.
    \item Tipos de Service e o modelo de rede do Kubernetes.
    \item Configuração, segredos e armazenamento persistente.
    \item Agendamento: requests, limits e afinidade.
    \item Alternativas de orquestração: Swarm e Nomad.
\end{itemize}
\end{frame}

\begin{frame}
\frametitle{Porque Precisamos de Orquestração}
\begin{itemize}
    \setlength\itemsep{0.4cm}
    \item O Docker, por si só, gere contentores num único host — não decide onde correr cada contentor, nem reage a falhas.
    \item Aplicações em produção precisam de vários hosts, para tolerância a falhas e para distribuir carga.
    \item É necessário decidir automaticamente em que nó correr cada carga de trabalho (agendamento/scheduling).
    \item É necessário detetar falhas e recuperar delas sem intervenção manual.
\end{itemize}
\end{frame}

\begin{frame}
\frametitle{Arquitetura Kubernetes: Control Plane \textcolor{NudeTaupe}{\small(1/2)}}
\begin{itemize}
    \setlength\itemsep{0.4cm}
    \item O control plane toma as decisões globais sobre o cluster (agendamento, deteção e resposta a eventos).
    \item \textbf{API server}: ponto de entrada único; todas as interações com o cluster passam por aqui.
    \item \textbf{etcd}: base de dados chave-valor distribuída que guarda o estado completo do cluster.
    \item \textbf{scheduler}: decide em que worker node cada novo Pod deve ser colocado.
\end{itemize}
\end{frame}

\begin{frame}
\frametitle{Arquitetura Kubernetes: Control Plane \textcolor{NudeTaupe}{\small(2/2)}}
\begin{itemize}
    \setlength\itemsep{0.4cm}
    \item \textbf{controller manager}: corre os controladores que comparam o estado atual com o estado desejado e agem para os aproximar.
    \item Os \textbf{worker nodes} são onde as cargas de trabalho efetivamente correm.
    \item \textbf{kubelet}: agente que garante que os contentores descritos para o nó estão de facto a correr.
    \item \textbf{kube-proxy} e o \textbf{container runtime} (ex.: containerd) completam o nó, tratando de rede e execução de contentores.
\end{itemize}
\end{frame}

\begin{frame}
\frametitle{Objetos Principais do Kubernetes \textcolor{NudeTaupe}{\small(1/2)}}
\begin{itemize}
    \setlength\itemsep{0.4cm}
    \item \textbf{Pod}: unidade mínima de execução; agrupa um ou mais contentores que partilham rede e armazenamento.
    \item \textbf{Deployment}: descreve o estado desejado de uma aplicação (imagem, número de réplicas, estratégia de atualização).
    \item O Deployment não gere Pods diretamente — delega isso num ReplicaSet.
\end{itemize}
\end{frame}

\begin{frame}
\frametitle{Objetos Principais do Kubernetes \textcolor{NudeTaupe}{\small(2/2)}}
\begin{itemize}
    \setlength\itemsep{0.4cm}
    \item \textbf{ReplicaSet}: garante que um determinado número de réplicas de um Pod está sempre em execução.
    \item \textbf{Service}: fornece um ponto de acesso estável (IP e nome) a um conjunto de Pods, que são efémeros por natureza.
    \item \textbf{Namespace}: divisão lógica do cluster, útil para isolar ambientes ou equipas dentro do mesmo cluster físico.
\end{itemize}
\end{frame}

\begin{frame}
\frametitle{Ciclo de Vida de um Pod e Self-Healing}
\begin{itemize}
    \setlength\itemsep{0.4cm}
    \item Um Pod passa por várias fases: Pending, Running, Succeeded, Failed.
    \item \textbf{Liveness probe}: verifica se o contentor continua saudável; se falhar, o kubelet reinicia-o.
    \item \textbf{Readiness probe}: verifica se o contentor está pronto para receber tráfego, antes de ser incluído num Service.
    \item Se um Pod morrer, o ReplicaSet cria automaticamente um novo, aproximando o cluster do estado desejado.
\end{itemize}
\end{frame}

\begin{frame}
\frametitle{Declarativo vs Imperativo}
\begin{itemize}
    \setlength\itemsep{0.4cm}
    \item Modo imperativo: comandos diretos que dizem ao Kubernetes exatamente \emph{o que fazer} (ex.: \texttt{kubectl run}, \texttt{kubectl create}).
    \item Modo declarativo: descreve-se o estado \emph{desejado} num manifesto YAML e o Kubernetes encarrega-se de o atingir.
    \item \texttt{kubectl apply -f manifesto.yaml} aplica (ou atualiza) o estado descrito no ficheiro.
    \item Abordagem declarativa é a recomendada em produção: reprodutível, versionável e auditável.
\end{itemize}
\end{frame}

\begin{frame}
\frametitle{O Ciclo de Reconciliação}
\begin{itemize}
    \setlength\itemsep{0.4cm}
    \item O Kubernetes funciona por \textbf{reconciliação contínua}: compara permanentemente o \emph{estado desejado} com o \emph{estado observado}.
    \item Cada controlador observa o seu tipo de recurso e atua para reduzir a diferença entre os dois estados.
    \item O estado desejado fica registado no \textbf{etcd}, através do API server; o estado observado é reportado pelos kubelets.
    \item Consequência prática: apagar um Pod gerido por um ReplicaSet não o elimina de forma duradoura -- o controlador recria-o.
    \item Este modelo, e não uma sequência de comandos, é o que dá ao Kubernetes a capacidade de self-healing.
\end{itemize}
\end{frame}

\begin{frame}
\frametitle{Tipos de Service}
\begin{itemize}
    \setlength\itemsep{0.4cm}
    \item \textbf{ClusterIP} (por omissão): expõe o serviço apenas dentro do cluster, num IP virtual interno.
    \item \textbf{NodePort}: abre a mesma porta em todos os nós do cluster, tornando o serviço acessível a partir do exterior.
    \item \textbf{LoadBalancer}: pede um balanceador externo ao fornecedor de cloud; em clusters on-premises requer uma solução adicional.
    \item \textbf{Ingress}: não é um tipo de Service, mas um recurso próprio que encaminha tráfego HTTP/HTTPS por nome de host ou caminho.
    \item A seleção dos Pods de destino é feita por \textbf{labels e selectors}, não por endereços fixos.
\end{itemize}
\end{frame}

\begin{frame}
\frametitle{O Modelo de Rede do Kubernetes}
\begin{itemize}
    \setlength\itemsep{0.4cm}
    \item Regra fundamental: cada Pod tem um endereço IP próprio e todos os Pods comunicam entre si sem NAT.
    \item Esta rede plana é implementada por um \textbf{plugin CNI} (Container Network Interface), não pelo Kubernetes em si.
    \item Exemplos de plugins CNI: Flannel, Calico, Weave Net, Cilium.
    \item \textbf{kube-proxy}: implementa em cada nó as regras que fazem um Service redirecionar tráfego para os Pods corretos.
    \item \textbf{CoreDNS}: fornece resolução de nomes interna, permitindo alcançar um Service pelo nome em vez do IP.
    \item Este modelo de rede será retomado na Aula 06, no contexto mais geral da SDN.
\end{itemize}
\end{frame}

\begin{frame}
\frametitle{Configuração e Segredos}
\begin{itemize}
    \setlength\itemsep{0.4cm}
    \item \textbf{ConfigMap}: guarda dados de configuração não sensíveis, separando-os da imagem do contentor.
    \item \textbf{Secret}: destinado a dados sensíveis (passwords, chaves, tokens).
    \item Atenção: por omissão, os Secrets são apenas codificados em base64 -- \emph{não} cifrados. É preciso ativar cifra em repouso no etcd e restringir acessos por RBAC.
    \item Ambos podem ser expostos ao contentor como variáveis de ambiente ou como ficheiros montados.
    \item Separar configuração da imagem permite usar a mesma imagem em desenvolvimento, teste e produção.
\end{itemize}
\end{frame}

\begin{frame}
\frametitle{Armazenamento Persistente}
\begin{itemize}
    \setlength\itemsep{0.4cm}
    \item Pods são efémeros: quando um Pod desaparece, o seu sistema de ficheiros local desaparece com ele.
    \item \textbf{PersistentVolume (PV)}: recurso de armazenamento disponibilizado no cluster.
    \item \textbf{PersistentVolumeClaim (PVC)}: pedido de armazenamento feito por uma aplicação, com capacidade e modo de acesso.
    \item \textbf{StorageClass}: permite aprovisionamento dinâmico, criando o volume automaticamente quando surge um PVC.
    \item \textbf{StatefulSet}: alternativa ao Deployment para aplicações com estado, que garante identidade e armazenamento estáveis por réplica.
    \item Estes conceitos ligam-se diretamente às Aulas 10 e 11, sobre armazenamento.
\end{itemize}
\end{frame}

\begin{frame}
\frametitle{Agendamento: Requests e Limits}
\begin{itemize}
    \setlength\itemsep{0.4cm}
    \item \textbf{Requests}: recursos que o Pod garantidamente precisa; o scheduler usa este valor para escolher um nó com capacidade suficiente.
    \item \textbf{Limits}: teto máximo que o contentor pode consumir, aplicado através dos cgroups vistos na Aula 04.
    \item Exceder o limite de memória leva à terminação do contentor (OOM kill); exceder o de CPU resulta apenas em \textit{throttling}.
    \item \textbf{nodeSelector} e \textbf{affinity}: influenciam em que nós um Pod pode ou deve ser colocado.
    \item \textbf{Taints e tolerations}: mecanismo inverso, em que um nó repele Pods que não tolerem explicitamente a sua marca.
\end{itemize}
\end{frame}

\begin{frame}
\frametitle{Alternativas de Orquestração}
\begin{itemize}
    \setlength\itemsep{0.4cm}
    \item \textbf{Docker Swarm}: solução de orquestração nativa do Docker, mais simples de configurar, mas com ecossistema e funcionalidades mais limitadas.
    \item \textbf{HashiCorp Nomad}: orquestrador genérico, capaz de agendar não só contentores mas também outras cargas de trabalho (binários, VMs).
    \item Kubernetes é hoje o standard de facto, com o ecossistema mais vasto e maior adoção na indústria.
    \item A escolha depende da complexidade do caso de uso e da curva de aprendizagem que a equipa está disposta a assumir.
\end{itemize}
\end{frame}

\begin{frame}
\frametitle{Síntese da Aula \textcolor{NudeTaupe}{\small(1/2)}}
\begin{itemize}
    \setlength\itemsep{0.4cm}
    \item Orquestração é necessária para gerir contentores de forma fiável em múltiplos hosts.
    \item Kubernetes separa control plane (API server, etcd, scheduler, controller manager) de worker nodes (kubelet, kube-proxy, runtime).
    \item Pod, Deployment, ReplicaSet, Service e Namespace são os objetos centrais para descrever e expor aplicações.
\end{itemize}
\end{frame}

\begin{frame}
\frametitle{Síntese da Aula \textcolor{NudeTaupe}{\small(2/2)}}
\begin{itemize}
    \setlength\itemsep{0.4cm}
    \item O ciclo de reconciliação contínua entre estado desejado e observado é o que dá ao Kubernetes a capacidade de self-healing.
    \item Services (ClusterIP, NodePort, LoadBalancer) e o modelo de rede plana via CNI expõem e interligam Pods efémeros.
    \item ConfigMaps, Secrets e PersistentVolumes separam configuração e dados do ciclo de vida dos Pods.
    \item Requests e limits guiam o agendamento e aplicam-se através dos cgroups vistos na Aula 04.
\end{itemize}
\end{frame}

\begin{frame}[plain]
  \vspace*{3cm}
  \begin{center}
    \Huge\bfseries\textcolor{TextEspresso}{Obrigada.}\par
    \vspace{0.5cm}
    {\color{NudeTaupe}\rule{2cm}{0.5pt}}\par
    \vspace{0.5cm}
    \Large\textcolor{TextEspresso}{\itshape Questões e comentários}
  \end{center}
\end{frame}

\end{document}
